\documentclass[conference]{IEEEtran}
\usepackage{cite}
\usepackage{graphicx}
\usepackage{booktabs}
\usepackage{url}

\title{Collaborative AI Agents on Low-Cost SBCs: An Evaluation in a Simulated Smart Home}

\author{\IEEEauthorblockN{Author list intentionally pending}}

\begin{document}
\maketitle

\begin{abstract}
This manuscript is a placeholder for an ongoing experiment. Results and conclusions will be written only after the protocol is frozen and the campaign data are collected.
\end{abstract}

\section{Introduction}
This study evaluates small AI agents running on low-cost Labrador SBCs in a simulated smart-home task. The title and claims are provisional.

\section{Research Question}
The primary question is whether collaboration among specialized agents improves task completion compared with a single agent, and what cost it adds in latency and processing.

\section{Method}
The experimental system uses real Labrador hardware and network communication, while the home, residents, sensors, appliances, and actions are simulated. A notebook coordinator maintains state, validates actions, and records structured logs. The coordinator uses explicit rules and does not use an additional language model.

\section{Protocol}
The planned comparison includes rule-based automation, one Labrador agent with full context, and three specialized Labrador agents. Development scenarios are used only for implementation and model selection. Reserved evaluation scenarios are used for the main campaign after protocol freeze.

\section{Results}
Results are intentionally pending.

\section{Limitations}
The experiment does not validate behavior in a real residence and does not measure real domestic energy consumption. Any simulated energy values are outputs of the simulator rather than measurements of appliances or Labrador boards.

\section{Conclusion}
Conclusions are intentionally pending.

\bibliographystyle{IEEEtran}
\bibliography{references}
\end{document}
